\chapter{Communication Complexity — Coin Tape}
%%% generated by the blueprint pipeline from TCSlib.CommunicationComplexity.NewmanTheorem.CoinTape
\section{Overview}

This module introduces the type of finite coin sequences used in randomised communication
complexity and equips it with the uniform probability measure.  A \emph{coin tape} of
length $n$ is a sequence of $n$ independent fair coin flips; the module provides the
canonical \texttt{MeasureSpace} instance whose volume is the uniform measure on all such
sequences, together with a proof that this measure is a probability measure.

%%%%%%%%%%%%%%%%%%%%%%%%%%%%%%%%%%%%%%%%%%%%%%%%%%%%%%%%%%%%%%%%%%%%%%%%%%%%%%%
\section{Declarations}

\begin{definition}[Coin tape of length $n$]
\lean{CommunicationComplexity.CoinTape}
\leanok
$\texttt{CoinTape}(n)$ is the type $\mathrm{Fin}\,n \to \mathrm{Bool}$, i.e.\ the set of
all binary strings of length $n$, representing the $2^n$ possible outcomes of $n$
independent fair coin flips.
\end{definition}
